\documentclass[11pt]{article}

\usepackage[margin=1in]{geometry}
\usepackage{amsmath}
\usepackage{amssymb}
\usepackage{booktabs}
\usepackage{enumitem}
\usepackage{hyperref}

% ---- notation macros -------------------------------------------------------
\newcommand{\Nindep}{N_{\mathrm{arr,indep}}}
\newcommand{\Nshare}{N_{\mathrm{arr,share}}}
\newcommand{\Ntot}{N_{\mathrm{arr,tot}}}
\newcommand{\Nsa}{N_{\mathrm{SA}}}
\newcommand{\Ndec}{N_{\mathrm{dec}}}
\newcommand{\Nbl}{N_{\mathrm{BL,arr}}}
\newcommand{\Nwl}{N_{\mathrm{WL,arr}}}
\newcommand{\tcyc}{t_{\mathrm{cycle}}}
\newcommand{\tdec}{t_{\mathrm{dec}}}
\newcommand{\twl}{t_{\mathrm{WL}}}
\newcommand{\tbl}{t_{\mathrm{BL}}}
\newcommand{\tsa}{t_{\mathrm{SA}}}
\newcommand{\tsw}{t_{\mathrm{switch}}}
\newcommand{\Volarr}{\mathrm{Vol}_{\mathrm{arr}}}
\newcommand{\Volsa}{\mathrm{Vol}_{\mathrm{SA}}}
\newcommand{\Voldec}{\mathrm{Vol}_{\mathrm{dec}}}
\newcommand{\Volsel}{\mathrm{Vol}_{\mathrm{sel}}}
\newcommand{\bacc}{b_{\mathrm{acc}}}          % bits/access
\newcommand{\Larr}{L_{\mathrm{arr}}}
\newcommand{\Lper}{L_{\mathrm{periph}}}
\newcommand{\Asa}{A_{\mathrm{SA}}}
\newcommand{\acell}{a_{\mathrm{cell}}}
\newcommand{\Volper}{\mathrm{Vol}_{\mathrm{periph}}}
\newcommand{\Eacc}{E_{\mathrm{access}}}
\newcommand{\Ebit}{E_{\mathrm{bit}}}
\newcommand{\Pdyn}{P_{\mathrm{dyn}}}
\newcommand{\Pmax}{P_{\max}}
\newcommand{\Aused}{A_{\mathrm{used}}}
\newcommand{\tlayer}{t_{\mathrm{layer}}}
\newcommand{\Ncell}{N_{\mathrm{cell,tot}}}

\title{A Macro-Level Bandwidth Model for Monolithic 3D Memory}
\author{3d\_memory \, --- design-space exploration tool}
\date{\today}

\begin{document}
\maketitle

\begin{abstract}
We give the analytical kernel behind \texttt{3d\_memory}, a macro-level
design-space-exploration tool for monolithic 3D memories at the
CACTI/DESTINY/NVSim altitude. Given a fixed storage capacity $C$, a fixed
footprint $A$, a fixed layer budget $L$ (each tier of physical thickness
$\tlayer$), and a peak power-density budget $\Pmax$ (a cooling limit set by an
external thermal stub), the tool searches 3D organizations that
\emph{maximize the internal array-supply bandwidth}. The central objects are a
throughput-bound cycle time, a bandwidth identity linear in the number of
independent peripheral sets, a volume-packing constraint that ties capacity,
cell density, and the layer partition together, and a per-access energy model
whose product with bandwidth is capped by $\Pmax$. Routing, the power-delivery
network, and the thermal boundary condition that produces $\Pmax$ are deferred
(Sec.~\ref{sec:scope}).
\end{abstract}

\section{Problem statement}
\label{sec:problem}

\begin{quote}
\emph{Given} a fixed capacity $C$, a fixed footprint $A$, a fixed number of
monolithic tiers $L$ (each of physical thickness $\tlayer$), and a peak
power-density budget $\Pmax$ (a cooling limit; the thermal stub of
Sec.~\ref{sec:scope} maps $\Delta T_{\max}$ onto it);
\emph{find} the 3D memory organization (array geometry, degree of peripheral
sharing, sense-amp stacking, and tier assignment of circuit classes) that
\emph{maximizes} the internal array-supply bandwidth $\mathrm{BW}$.
\end{quote}

Two named 3D levers act on this objective:
\begin{enumerate}[nosep]
  \item \textbf{Peripheral stacking} --- tucking sense amps, decoders, and
        selectors on tiers under/over the array (CUA / Choi-style), freeing
        footprint and relieving pitch.
  \item \textbf{Multi-tier bitcells} --- placing different devices of a single
        cell on different tiers, raising areal cell density.
\end{enumerate}

\section{Nomenclature}
\label{sec:notation}

\begin{center}
\begin{tabular}{@{}l p{0.70\textwidth}@{}}
\toprule
Symbol & Meaning \\
\midrule
$C,\;A,\;L$        & storage capacity, footprint (area/layer), layer budget \\
$\tlayer$          & physical thickness per tier (so $\mathrm{VolBudget}=A\,L\,\tlayer$) \\
$\Delta T_{\max}$  & peak temperature rise allowed (thermal governor, feeds $\Pmax$) \\
$\Pmax$            & peak power-density budget [$\mathrm{W/mm^2}$] (cooling limit) \\
$\acell$           & areal footprint of one stored bit (cell density$^{-1}$) \\
$\Nindep$          & number of \emph{independent} peripheral sets (arrays or blocks of shared arrays) \\
$\Nshare$          & number of arrays time-multiplexed onto one shared peripheral set \\
$\Ntot$            & total array count, $\Ntot=\Nindep\cdot\Nshare$ \\
$\Nsa,\;\Ndec$     & number of sense amps, number of decoders \\
$\bacc$            & bits delivered per access (sense width) \\
$\Nbl,\;\Nwl$      & bitlines, wordlines per array \\
$\tcyc$            & steady-state cycle time (throughput bound) \\
$\tdec,\twl,\tbl,\tsa$ & decode, wordline, bitline-develop, and sense (incl.\ restore) times \\
$\tsw$             & selector switching time for time-muxing shared arrays \\
$\Volarr,\Volsa,\Voldec$ & volumes of one array, sense amp, decoder \\
$\Volsel,\Volper$  & volumes of selector network, fixed periphery per array \\
$\Larr,\Lper$      & layers holding bits vs.\ layers holding periphery \\
$\Asa$             & footprint of one sense amp (constant) \\
$K$                & column-multiplexing degree \\
$w_{\mathrm{SA}},\,p_{\mathrm{BL}}$ & sense-amp pitch, bitline pitch \\
$u$                & areal utilization of a peripheral tier \\
$\Eacc,\Ebit$      & per-access energy, per-delivered-bit energy \\
$\Pdyn$            & dynamic power ($=\mathrm{BW}\cdot\Ebit$) \\
$\Aused$           & footprint the packed volume actually occupies \\
$\Ncell$           & total stored cells $=\Nbl\Nwl\Ntot$ \\
\bottomrule
\end{tabular}
\end{center}

\section{The kernel}
\label{sec:kernel}

\subsection{Cycle time (throughput bound)}
The steady-state cycle time is set by the slowest \emph{un-hideable} resource
occupancy. The shared decoder and sense-amp occupancies are floors that cannot
be pipelined away; the array develop latency $\twl+\tbl$ can be hidden by
interleaving $\Nshare$ arrays onto the shared periphery:
\begin{equation}
  \tcyc = \max\!\left(\;
    \tdec,\;\;
    \tsa+\tsw,\;\;
    \frac{\tdec+\twl+\tbl+\tsa+\tsw}{\Nshare}
  \;\right).
  \label{eq:tcyc}
\end{equation}
$\tdec$ is the critical decoder-stage delay; $\tsa$ \emph{includes restore} for
destructive-read cells. This last inclusion silently encodes the technology
dependence of sharing: for a destructive cell (DRAM) $\tsa$ floors near $t_{RC}$
and the shared term can never dominate, so interleaving buys nothing; for a
non-destructive cell (gain cell, NVM) $\tsa$ is small and interleaving pays.

\subsection{Bandwidth}
Internal array-supply bandwidth is the per-access word width divided by the
cycle time, replicated across every \emph{independent} peripheral set:
\begin{equation}
  \boxed{\;\mathrm{BW} = \frac{\bacc}{\tcyc}\,\Nindep\;}
  \qquad\text{with}\qquad
  \bacc \le \Nbl .
  \label{eq:bw}
\end{equation}
The cap $\bacc\le\Nbl$ states that no access can sense more bits than there are
bitlines (whole-row-activate limit). Sharing raises $\Ntot$ and hides latency
but never appears in \eqref{eq:bw}: it adds capacity, not bandwidth.

\subsection{Accounting identities}
\begin{align}
  \Ntot &= \Nindep\cdot\Nshare, &
  \Nsa  &= \Nindep\cdot\bacc, &
  \Ndec &= \Nindep .
  \label{eq:counts}
\end{align}
There is one decoder per independent set, and the sense-amp count is the total
delivered word width.

\subsection{Volume (packing identity)}
Consumed volume equals the sum of every packed component --- an equality, not a
bound, so that no volume goes unaccounted:
\begin{equation}
  \mathrm{Vol} \;=\;
    \Volarr\cdot\Ntot \;+\;
    \Volsa\cdot\Nsa   \;+\;
    \Voldec\cdot\Ndec \;+\;
    \Volsel(\Nshare) \;+\;
    \Volper\cdot\Ntot .
  \label{eq:vol}
\end{equation}
The selector term $\Volsel(\Nshare)$ is the volume of the network that
time-multiplexes $\Nshare$ arrays onto one shared peripheral set; it grows with
the sharing degree (concretely $\Volsel\propto(\Nshare-1)\,\Nsa$, so it
vanishes at $\Nshare=1$) and is what makes sharing cost area, not just time. The
final term $\Volper\cdot\Ntot$ is a fixed per-array peripheral overhead (row
decoders, wordline drivers, array edge). This same polynomial is reused as the
occupied footprint $\Aused=\mathrm{Vol}/(L\,\tlayer)$ in the power-density
constraint below.

\subsection{Energy and power density}
Each single-array read activates a whole $\Nbl$-wide row but delivers only
$\bacc$ bits (O'Connor row overfetch). The per-access dynamic energy is a
uniform $CV^2$ form plus fixed periphery and an optional write term:
\begin{equation}
  \Eacc \;=\; k_{\mathrm{col}}\,\Nbl \;+\; k_{\mathrm{arr}}\,\Nbl\Nwl
         \;+\; e_{\mathrm{periph}} \;+\; f_w\,e_{\mathrm{write}}\,\bacc,
  \label{eq:eacc}
\end{equation}
with $k_{\mathrm{col}}=C_{\mathrm{cell}}V^2$ (the per-bitline cell/access cap)
and $k_{\mathrm{arr}}=c_{BL}V^2$ (the activated row's distributed bitline cap,
growing with rows). The swing is $V^2=V_{\mathrm{read}}^2$ for full-rail modes
(SRAM latch, resistive-read precharge, and destructive charge-share write-back);
the partial product $V_{\mathrm{read}}V_{\mathrm{sense}}$ applies only to a
hypothetical \emph{non-destructive} charge read. Amortizing over the delivered
bits and multiplying by bandwidth gives the dynamic power, which together with a
flat per-bit leakage/refresh term is held under the cooling budget as a
\emph{power density}:
\begin{align}
  \Ebit &= \frac{\Eacc}{\bacc}, &
  \Pdyn &= \mathrm{BW}\cdot\Ebit, &
  \frac{\Pdyn + p_{\mathrm{leak}}\,\Ncell}{\Aused} &\le \Pmax,
  \label{eq:power}
\end{align}
with $\Aused=\mathrm{Vol}/(L\,\tlayer)$ from \eqref{eq:vol}. In the solver the
density constraint is written division-free (moving $\Aused$, itself the volume
polynomial, to the right-hand side) so it stays polynomial for the spatial
branch-and-bound. Because the footprint reuses the packed volume, energy,
volume, and geometry are coupled through one expression; $\Pmax$ is supplied by
the external thermal stub of Sec.~\ref{sec:scope}.

\subsection{Scaling laws}
\begin{align}
  \twl &= k^{\mathrm{wire}}_{WL}\,\Nbl^2 + k^{\mathrm{cell}}_{WL}\,\Nbl, &
  \tbl &= f_\Delta\,\bigl(a_{\mathrm{lin}}\,\Nwl + a_{\mathrm{quad}}\,\Nwl^2\bigr), \\
  \Voldec &\propto \Nwl\cdot\Nbl \propto \Volarr, &
  \Volsa  &= \text{const}.
  \label{eq:scaling}
\end{align}
The wordline develop $\twl$ keeps the wire/cell split: a distributed-line
self-RC term $\tfrac12 rc L^2$ (CACTI-6 $n(n{+}1)rc/2$ form, quadratic in the
$\Nbl$ cells along the line; coefficient $k^{\mathrm{wire}}_{WL}=\tfrac12 rc$
set by the BEOL stack, \emph{shared} across cell technologies) plus a linear
gate-loading term $k^{\mathrm{cell}}_{WL}$ (\emph{per technology}).

The bitline develop $\tbl$ is the \textbf{sense-signal settling time}, derived
per sensing mode from Upton (2024) (current-mode readout, App.~A; voltage-mode
readout, App.~B), with a charge-sharing branch for destructive DRAM/eDRAM cells.
All three reduce to the same polynomial in the $\Nwl$ cells on the bitline: a
linear \emph{signal-development} term $a_{\mathrm{lin}}$ plus a quadratic BL wire
self-RC term $a_{\mathrm{quad}}=\tfrac12 r_{BL}c_{BL}$, scaled by a settling
factor $f_\Delta=-\ln(1-\Delta)$ that carries the sensing margin as a \emph{time}
cost. The settling fraction $\Delta$ (e.g.\ $0.99$) is a fixed reliability
requirement, not a free variable; tightening it costs time only
logarithmically. The linear term is mode-specific:
\begin{itemize}[nosep]
  \item \textbf{current} (App.~A, Eq.~A.12): development time constant
        $\tau_C=C_{BL}\,nU_t/I_{\mathrm{read}}$, so
        $a_{\mathrm{lin}}=c_{BL}\,nU_t/I_{\mathrm{read}}$.
  \item \textbf{voltage} (App.~B, Eq.~B.7): Elmore constant
        $\tau=v_{\mathrm{ratio}}(R_{\text{pull-up}}+R_{BL}/2)C_{BL}$,
        pull-up-dominated, giving
        $a_{\mathrm{lin}}=v_{\mathrm{ratio}}R_{\text{pull-up}}c_{BL}$; the same
        divider ratio $v_{\mathrm{ratio}}=V_{BL,\mathrm{cell}}/V_{\mathrm{read}}$
        also multiplies the wire term $a_{\mathrm{quad}}$. We use the single-$\tau$
        settling latency $f_\Delta=-\ln(1-\Delta)$, \emph{not} the dual-edge
        clock-period bound $T_{\mathrm{clk}}\ge-2\tau\ln(1-\Delta)$ of Eq.~B.8
        --- the $2\times$ is a half-cycle convention on the clock \emph{period},
        not on the serial develop latency, and dropping it keeps the shared wire
        term consistent across modes.
  \item \textbf{charge\_share} (DRAM/eDRAM, DESTINY convention): passive charge
        redistribution with a lumped series-RC
        $a_{\mathrm{lin}}=r_{BL}C_{\mathrm{cell}}$ (CACTI-D / CACTI-5.1) plus the
        same distributed wire self-RC $a_{\mathrm{quad}}$, which overtakes the
        lumped term once $\Nwl\gtrsim 2C_{\mathrm{cell}}/c_{BL}$ --- the long,
        thin, resistive 3D-BEOL bitline regime this tool targets. The
        series-cap ratio governs signal \emph{amplitude}, not speed, so it
        enters only as the row cap of Sec.~\ref{sec:derived}.
\end{itemize}
$\tsa$ retains the intrinsic device read/latch floor (and restore, for
destructive cells). Decoder volume tracks the array it drives; sense-amp volume
is a fixed footprint per amp --- the earlier Pelgrom area law
$\Asa\propto1/\mathrm{margin}^2$ is dropped, because the accurate cost of a
sensing margin is the settling \emph{time} above, not sense-amp area.

\section{Derived structure}
\label{sec:derived}

\subsection{Layer partition}
Capacity and footprint fix how many tiers must hold bits; the rest are free for
periphery:
\begin{align}
  \Larr  &= \left\lceil \frac{C\,\acell}{A} \right\rceil, &
  \Lper  &= L - \Larr, &
  \Nsa   &\le \frac{A\,\Lper\,u}{\Asa}.
  \label{eq:layers}
\end{align}
Bandwidth lives in $\Lper$. Hence \textbf{cell density is a quantified indirect
bandwidth lever}: shrinking $\acell$ (e.g.\ by making the cell 3D / multi-tier)
lowers $\Larr$ and frees peripheral tiers for more sense amps --- worthwhile iff
the freed tiers buy more $\Nsa$ than the added cell complexity costs in $\Volarr$.

\subsection{Charge-share row cap}
For a destructive charge-sharing cell the developed signal amplitude
$C_{\mathrm{cell}}/(C_{\mathrm{cell}}+c_{BL}\Nwl)$ must stay above the sense-amp
offset margin $m_{\mathrm{SA}}$. This is an amplitude, not a speed, constraint,
so it enters as a constant upper bound on the rows rather than a develop-time
term:
\begin{equation}
  \Nwl \;\le\; \frac{C_{\mathrm{cell}}\,(1-m_{\mathrm{SA}})}{c_{BL}\,m_{\mathrm{SA}}}.
  \label{eq:rowcap}
\end{equation}
This is why DESTINY caps DRAM subarray rows; for non-destructive
(current/voltage) cells the bound is absent.

\subsection{Optimal sharing degree}
Interleaving helps only until the develop term in \eqref{eq:tcyc} drops to the
floor; beyond that, more sharing adds $\Volsel$ and $\Ntot$ for no bandwidth:
\begin{equation}
  \Nshare^\star =
  \left\lceil
    \frac{\tdec+\twl+\tbl+\tsa+\tsw}{\max(\tdec,\;\tsa+\tsw)}
  \right\rceil .
  \label{eq:nshare}
\end{equation}

\subsection{Width / mux / pitch coupling}
The sense width is set by the column-mux degree, which is floored by the
sense-amp pitch; stacking $n$ sense-amp tiers relieves that pitch:
\begin{equation}
  \bacc = \frac{\Nbl}{K},
  \qquad
  K \ge \left\lceil \frac{w_{\mathrm{SA}}}{p_{\mathrm{BL}}} \right\rceil \Big/\, n .
  \label{eq:width}
\end{equation}
So the width lever, the mux degree, the sense-amp pitch, and sense-amp tier
stacking are a single coupled knob.

\section{Allocation recipe}
\label{sec:recipe}
Given $(C,A,L,\text{cell/tier tech})$, the optimizer walks the levers in
dependency order:
\begin{enumerate}[nosep]
  \item \textbf{Cell density} --- fixes $\Larr$ via \eqref{eq:layers}, freeing $\Lper$.
  \item \textbf{Array geometry} --- fixes the $\tcyc$ floor \eqref{eq:tcyc} and $\Volarr,\Asa$.
  \item \textbf{Sense-amp stacking + mux degree} --- raises $\bacc$ toward $\Nbl$ \eqref{eq:width}.
  \item \textbf{Share to $\Nshare^\star$} --- reaches the $\tcyc$ floor at minimum cost \eqref{eq:nshare}.
  \item \textbf{Fill $\Lper$} --- spend remaining peripheral area on $\Nsa$ \eqref{eq:layers}, maximizing \eqref{eq:bw}, until either the peripheral area or the power-density budget \eqref{eq:power} binds.
\end{enumerate}

\section{Scope and deferrals}
\label{sec:scope}
The model is a \emph{packing-plus-power} model: how much bandwidth-producing
hardware fits in $L$ tiers of area $A$ at capacity $C$ without exceeding the
cooling power density $\Pmax$ \eqref{eq:power}. Two assumptions justify deferring
routing:
(i) \textbf{distributed vertical bottom-exit} --- each independent array exits
its bits through a local per-array vertical port grid, so there is no die-scale
horizontal data-collection network; and (ii) \textbf{consumer-provisioned} --- we
model that the bandwidth exists at the exit surface, not who consumes it. Under
these, data egress collapses to a short local vertical hop (roughly uniform
across banking choices), so the packing result is close to physical rather than a
loose upper bound. The internal selector cost is retained ($\tsw$, $\Volsel$).

\medskip\noindent
\textbf{Deferred:} the thermal boundary condition that \emph{sets} $\Pmax$ (a
pluggable stub maps an areal $R_{\mathrm{th}}$ and $\Delta T_{\max}$ to the
power-density budget consumed in \eqref{eq:power}); the power-delivery network;
and the lateral address/command distribution (narrow, assumed provisioned). Note
the power-density \emph{constraint} itself is no longer deferred --- it is
enforced directly through the per-access energy model of Sec.~\ref{sec:kernel}.

\end{document}
